\documentclass[12pt,a4paper]{article}
\usepackage[utf8]{inputenc}
\usepackage[portuguese]{babel}
\usepackage{geometry}
\usepackage{listings}
\usepackage{xcolor}
\usepackage{hyperref}
\usepackage{geometry}

\geometry{margin=1in}
\setlength{\parindent}{0pt}

\lstset{
    keywordstyle=\color{blue},
    commentstyle=\color{gray},
    stringstyle=\color{green},
    frame=single,
    tabsize=2,
    showstringspaces=true,
    breaklines=true,
    basicstyle=\ttfamily\small,
    columns=fixed
}

\begin{document}

\begin{titlepages}

\centering

\vspace*{1cm}

{\LARGE INSTITUTO FEDERAL DO ESPÍRITO SANTO \par}

\vspace{0.5cm}

{\large CAMPUS CACHOEIRO DE ITAPEMIRIM \par}

\vspace{1.5cm}

{\Large RELATÓRIO TÉCNICO \par}

\vspace{1.5cm}

{\Large \textbf{Sistemas Distribuídos Heterogêneos: Integração entre Cliente JavaFX e Servidor gRPC em Python} \par}

\vspace{2cm}

{\large Autores: \par}
\vspace{0.3cm}
{\large André Mendonça Valane \\ Marcos Lopes Ribeiro \par}

\vspace{2cm}

{\large Disciplina: Sistemas Distribuídos \par}
\vspace{0.3cm}
{\large Professor: João Paulo de Brito Gonçalves \par}
\vspace{0.3cm}
{\large Curso: Bacharelado em Sistemas de Informação \par}

\vspace{2cm}

{\large Cachoeiro de Itapemirim, \today \par}

\end{titlepages}

\newpage

\tableofcontents

\newpage

\section{Introdução}

Os sistemas distribuídos heterogêneos representam um desafio significativo na arquitetura de software moderna, exigindo a integração de componentes desenvolvidos em linguagens e paradigmas diferentes. Neste trabalho, exploramos a integração de uma aplicação cliente construída com JavaFX — uma tecnologia orientada a objetos voltada para interfaces gráficas — com um servidor gRPC escrito em Python. A escolha do gRPC como protocolo de comunicação foi motivada pela sua eficiência em transferência de mensagens e suporte nativo a tipagem forte através de Protocol Buffers (Protobuf).

O objetivo principal deste projeto foi implementar uma funcionalidade de consulta remota de laboratórios por ID, onde o cliente envia um identificador e o servidor retorna as informações detalhadas do laboratório. Essa tarefa demonstra a interoperabilidade entre ecossistemas distintos: JavaFX (desenvolvimento de UI) e Python (servidor de serviços), unificados por um contrato definido em Protocol Buffers.

A heterogeneidade das tecnologias foi tratada através de um servidor gRPC que expõe operações via RPCs, enquanto o cliente JavaFX utiliza threads secundárias para realizar as chamadas de rede, garantindo que a interface gráfica permaneça responsiva durante a comunicação remota. O banco de dados PostgreSQL serve como fonte de verdade compartilhada entre ambos os componentes.

\section{Arquitetura e o Contrato gRPC}

A comunicação entre cliente e servidor é mediada por um contrato definido em Protocol Buffers, descrito no arquivo \texttt{laboratorio.proto}. Este arquivo define as mensagens principais utilizadas na comunicação remota.

\begin{lstlisting}[language=protobuf, caption={\'Contrato Protocol Buffers}]
syntax = "proto3";

package laboratorio;

// Mensagem para solicitar informações de laboratório
message LaboratorioRequest {
    int64 id = 1;
}

// Mensagem para responder com informações de laboratório
message LaboratorioResponse {
    int64 id = 1;
    string nome = 2;
    string descricao = 3;
    string localizacao = 4;
    int64 numero_horas = 5;
}

// Serviço gRPC que expõe a operação de consulta
service LaboratorioService {
    // RPC para buscar laboratório por ID
    rpc BuscarLaboratorio (LaboratorioRequest) returns (LaboratorioResponse);
}
\end{lstlisting}

O contrato utiliza campos numéricos para identificadores e strings para metadados, garantindo clareza e compatibilidade entre as partes da comunicação. As mensagens são serializadas automaticamente pelo Protobuf, permitindo compressão e transferência eficiente sobre redes TCP/IP.

O servidor gRPC é implementado em Python usando a biblioteca \texttt{grpcio} e \texttt{grpcio-tools}, que facilitam a geração de código estrutural a partir do arquivo .proto. O cliente, por sua vez, consome este mesmo contrato para enviar requisições e receber respostas.

\section{Implementação do Servidor (Python)}

O servidor foi implementado em Python, utilizando a biblioteca \texttt{grpcio} para definir e executar o serviço gRPC, e \texttt{psycopg2} para interagir com o banco de dados PostgreSQL.

A estrutura principal do servidor (\texttt{server.py}) inclui:
\begin{itemize}
    \item Definição do serviço \texttt{LaboratorioService} herdado do gRPC
    \item Implementação do método \texttt{BuscarLaboratorio} que recebe um \texttt{LaboratorioRequest}
    \item Consulta parametrizada ao PostgreSQL para obter os dados do laboratório
    \item Retorno da resposta formatada conforme o contrato \texttt{LaboratorioResponse}
    \item Tratamento adequado de códigos de status gRPC, incluindo \texttt{NOT\_FOUND} quando o laboratório não existe e \texttt{INTERNAL} para erros internos do servidor
\end{itemize}

\begin{lstlisting}[language=python, caption={\'Servidor gRPC em Python}]
#!/usr/bin/env python3

import grpc
from concurrent import futures
from google.rpc import code_pb2, status_pb2
from google.rpc.status_pb2 import Status
import psycopg2
from laboratorio_pb2 import LaboratorioRequest, LaboratorioResponse
from laboratorio_pb2_grpc import LaboratorioServiceServicer, add_LaboratorioServiceServicer_to_server

class LaboratorioServicer(LaboratorioServiceServicer):
    def __init__(self):
        self.conn = None

    def connect_db(self):
        try:
            self.conn = psycopg2.connect(
                host="localhost",
                database="laboratorios",
                user="postgres",
                password="senha_do_banco"
            )
        except Exception as e:
            raise Exception(f"Erro ao conectar ao banco de dados: {e}")

    def BuscarLaboratorio(self, request, context):
        if not self.conn:
            self.connect_db()

        try:
            cursor = self.conn.cursor()
            cursor.execute(
                "SELECT id, nome, descricao, localizacao, numero_horas "
                "FROM laboratorios WHERE id = %s",
                (request.id,)
            )
            result = cursor.fetchone()
            cursor.close()

            if result is None:
                context.set_code(grpc.StatusCode.NOT_FOUND)
                context.set_details(f"Laboratório com ID {request.id} não encontrado")
                return LaboratorioResponse()

            response = LaboratorioResponse(
                id=result[0],
                nome=result[1],
                descricao=result[2],
                localizacao=result[3],
                numero_horas=result[4]
            )
            return response

        except Exception as e:
            context.set_code(grpc.StatusCode.INTERNAL)
            context.set_details(f"Erro interno do servidor: {e}")
            return LaboratorioResponse()

def serve():
    server = grpc.server(futures.ThreadPoolExecutor(max_workers=10))
    add_LaboratorioServiceServicer_to_server(LaboratorioServicer(), server)
    server.add_insecure_port('[::]:50051')
    server.start()
    print("Servidor gRPC iniciado na porta 50051")
    server.wait_for_termination()

if __name__ == '__main__':
    serve()
\end{lstlisting}

O servidor utiliza \texttt{ThreadPoolExecutor} para escalar o processamento de múltiplas requisições simultaneamente, garantindo alta disponibilidade. A conexão com o PostgreSQL é estabelecida no momento da inicialização do servicer, e consultas são realizadas de forma segura através de parâmetros, evitando injeção de SQL.

\section{Implementação do Cliente (JavaFX)}

O cliente foi desenvolvido em JavaFX, aproveitando o sistema de laboratórios de POO II para estruturação modular. A principal desafio foi garantir que as chamadas de rede não travassem a interface gráfica, o que foi resolvido isolando essas operações em threads secundárias.

O cliente utiliza o plugin Maven para gerenciamento de dependências e construção do projeto. A automação das chamadas gRPC é realizada através de um cliente gRPC customizado que encapsula o contrato \texttt{laboratorio.proto}.

\begin{lstlisting}[language=java, caption={\'Cliente JavaFX com gRPC}]
package com.example.laboratorio;

import javafx.application.Application;
import javafx.application.Platform;
import javafx.geometry.Insets;
import javafx.geometry.Pos;
import javafx.scene.Scene;
import javafx.scene.control.Alert;
import javafx.scene.control.Button;
import javafx.scene.control.Label;
import javafx.scene.control.TextField;
import javafx.scene.layout.HBox;
import javafx.scene.layout.VBox;
import javafx.stage.Stage;
import io.grpc.ManagedChannel;
import io.grpc.ManagedChannelBuilder;
import laboratorio.LaboratorioRequest;
import laboratorio.LaboratorioResponse;
import laboratorio.LaboratorioServiceGrpc;

public class LaboratorioClient extends Application {

    private TextField idField;
    private TextField nomeField;
    private TextField descricaoField;
    private TextField localizacaoField;
    private TextField horasField;
    private ManagedChannel channel;

    @Override
    public void start(Stage primaryStage) {
        channel = ManagedChannelBuilder
            .forAddress("localhost", 50051)
            .usePlaintext()
            .build();

        idField = new TextField();
        nomeField = new TextField();
        descricaoField = new TextField();
        localizacaoField = new TextField();
        horasField = new TextField();

        idField.setPromptText("ID do Laboratório");
        nomeField.setEditable(false);
        descricaoField.setEditable(false);
        localizacaoField.setEditable(false);
        horasField.setEditable(false);

        Button searchButton = new Button("Consultar");
        searchButton.setOnAction(e -> buscarLaboratorio());

        HBox inputBox = new HBox(10, idField, searchButton);
        inputBox.setAlignment(Pos.CENTER_LEFT);
        inputBox.setPadding(new Insets(10));

        VBox formBox = new VBox(10);
        formBox.getChildren().addAll(
            new Label("Nome:"), nomeField,
            new Label("Descrição:"), descricaoField,
            new Label("Localização:"), localizacaoField,
            new Label("Número de Horas:"), horasField
        );
        formBox.setPadding(new Insets(10));

        VBox root = new VBox(10, inputBox, formBox);
        root.setPadding(new Insets(20));

        Scene scene = new Scene(root, 500, 400);
        primaryStage.setTitle("Consulta de Laboratórios - Sistemas Distribuídos");
        primaryStage.setScene(scene);
        primaryStage.show();
    }

    private void buscarLaboratorio() {
        String idText = idField.getText();
        if (idText.isEmpty()) {
            showAlert("Por favor, insira um ID válido.");
            return;
        }

        long id;
        try {
            id = Long.parseLong(idText);
        } catch (NumberFormatException ex) {
            showAlert("ID inválido. Insira um número inteiro.");
            return;
        }

        // Chamada assíncrona para não bloquear a UI
        new Thread(() -> {
            try {
                LaboratorioServiceGrpc.LaboratorioServiceBlockingStub stub =
                    LaboratorioServiceGrpc.newBlockingStub(channel);

                LaboratorioRequest request = LaboratorioRequest.newBuilder()
                    .setId(id)
                    .build();

                LaboratorioResponse response = stub.buscarLaboratorio(request);

                // Uso de Platform.runLater para atualizar a UI no thread principal
                Platform.runLater(() -> {
                    nomeField.setText(response.getNome());
                    descricaoField.setText(response.getDescricao());
                    localizacaoField.setText(response.getLocalizacao());
                    horasField.setText(String.valueOf(response.getNumeroHoras()));
                });

            } catch (Exception ex) {
                Platform.runLater(() -> {
                    showAlert("Erro ao consultar laboratório: " + ex.getMessage());
                });
            }
        }).start();
    }

    private void showAlert(String message) {
        Alert alert = new Alert(Alert.AlertType.ERROR);
        alert.setTitle("Erro");
        alert.setHeaderText(null);
        alert.setContentText(message);
        alert.showAndWait();
    }

    @Override
    public void stop() {
        if (channel != null && !channel.isShutdown()) {
            channel.shutdownNow();
        }
    }

    public static void main(String[] args) {
        launch(args);
    }
}
\end{lstlisting}

O cliente segue o princípio de separação de preocupações: a lógica de negócios (consulta ao servidor remoto) está separada da camada de interface gráfica. O uso de \texttt{Thread} garante que a thread principal da JavaFX permaneça livre para atualizar a interface, enquanto o processamento de rede ocorre em segundo plano. O método \texttt{Platform.runLater} é crucial para garantir que qualquer modificação na interface seja feita de forma segura, pois ele agenda a execução no thread principal da JavaFX.

As chamadas gRPC são enviadas de forma síncrona dentro da thread secundária, e o resultado é exibido na interface gráfica através do \texttt{Platform.runLater}, que garante a atualização segura da UI no thread principal.

\section{Considerações Finais}

O trabalho apresentado demonstrou a viabilidade da integração entre tecnologias heterogêneas — JavaFX e Python — através de um contrato gRPC bem definido. Os resultados confirmaram que o gRPC oferece vantagens significativas em relação a abordagens legadas baseadas em sockets tradicionais:

\begin{enumerate}
    \item \textbf{Eficiência de Transferência}: O gRPC utiliza protocolos binários compactados, reduzindo o tamanho das mensagens em comparação com sockets TCP puros.
    \item \textbf{Tipagem Forte}: O uso de Protocol Buffers garante que os contratos entre cliente e servidor sejam verificados em tempo de compilação, eliminando erros de integração.
    \item \textbf{Desempenho}: A comunicação gRPC é otimizada para baixa latência e alto throughput, ideal para aplicações de sistemas distribuídos.
    \item \textbf{Manutenibilidade}: O contrato definido em Protobuf facilita a evolução futura do sistema, permitindo adição de novos campos sem quebrar clientes existentes.
\end{enumerate}

A escolha do Python para o servidor foi motivada pela maturidade da biblioteca gRPC e pela facilidade de integração com bancos de dados relacionais como o PostgreSQL. Já o JavaFX foi selecionado para o cliente devido à sua popularidade no desenvolvimento de interfaces gráficas e à capacidade de integrar facilmente com bibliotecas Java.

Em suma, esta implementação representa um exemplo prático de sistemas distribuídos heterogêneos, onde a interoperabilidade é alcançada através de padrões modernos de comunicação e armazenamento de dados.

\begin{thebibliography}{9}
\bibitem{google2018grpc}
Google. \textit{gRPC: Remote Procedure Calls for the Cloud} (2018). Available at: \url{https://cloud.google.com/grpc/docs/language-guide/overview}

\bibitem{protobuf}
Google. \textit{Protocol Buffers: A Language Neutral Interface Definition Language} (2015). Available at: \url{https://developers.google.com/protocol-buffers/docs/introduction}

\bibitem{postgresql}
PostgreSQL Global Development Group. \textit{PostgreSQL Documentation}. Available at: \url{https://www.postgresql.org/docs/}

\bibitem{javafx}
OpenJFX. \textit{JavaFX: A Modern UI Framework for Java Applications}. Available at: \url{https://openjfx.io/}

\bibitem{grpcio}
gRPC Python. \textit{gRPC Implementation in Python}. Available at: \url{https://github.com/grpc/grpcio}

\end{thebibliography}

\end{document}
